\documentclass[
  spinewidth=12mm,
  coverheight=190mm,
  coverwidth=125mm,
  bleedwidth=3mm,
  marklength=0mm
]{bookcover}

\usepackage[T1]{fontenc}
\usepackage[sfdefault]{AlegreyaSans}
\usepackage{xcolor}

\definecolor{stars}{HTML}{02174c}
\newcommand{\booktitle}{Das Gastmahl der Schatten}
\newcommand{\bookauthor}{Julius Siber}
\newcommand{\booksubtitle}{Roman
von einer\\ polnischen Allerseelennacht\\
und Friedrich Chopin}

\newbookcovercomponenttype{center rotate}{%
  \parbox[t][\partheight][c]{\partwidth}{%
    \begin{center}
       \rotatebox[origin=c]{90}{#1}
    \end{center}}}

\begin{document}
  \begin{bookcover}
    \bookcovercomponent{color}{bg whole}{black!90!white}
    \bookcovercomponent{color}{bg spine}{black!90!white}
    \bookcovercomponent{color}{bg front}{stars!70!black}
    \bookcovercomponent{normal}{front}[0.06\textwidth,0.06\textwidth,0.06\textwidth,0.06\textwidth]{
      \Huge\color{white}
      {\Large\bookauthor{}\\[0.5\baselineskip]}
      {\Huge\bfseries Das Gastmahl \\ der Schatten\\[0.5\baselineskip]}
      {\huge\booksubtitle{}}\vfill
      \par
      PTOROS}


  \bookcovercomponent{normal}{back}{%
    \centering\large\color{white}\vspace{2em}
    \parbox{0.38\textwidth}{%
      \textit{Das Gastmahl der Schatten}, letzter, rätselhaftester und verschollenster Roman des zu Lebzeiten weltberühmten Seelenfängers und Geigenvirtuosen \textbf{Julius Siber} (1871–1943), erschien 1937 als wortgewaltige Wiederheraufbeschwörung und zugleich literarische Letztverwahrung jenes okkultistisch durchwehten Kunstfanatismus, der über weite Strecken des 19. Jahrhunderts die kulturellen Zentren und bedeutendsten Geister Europas in Bann schlug. Folgte die 1914 veröffentlichte Seelenwanderung, Sibers heute wohl bekanntester Roman, noch den Spuren eines Malers, dessen übernatürliche Ambition sich keineswegs in der Arbeit an der Leinwand erschöpft – denn niemals wird bei Siber bloß die Kunst durch ihren Schöpfer zur Magie, immer wächst vor allem ihr Schöpfer an der Kunst erst recht zum Magier empor –, so kehrt der fränkische Schriftsteller und Komponist sechs Jahre vor seinem Tod triumphal-mitreißend noch einmal zur eigenen und eigentlichsten Dämonie-Domäne und damit zu jener Kunstgattung zurück, der seit jeher die größte Zwingkraft über Sinne und Herzen der Menschen zugeschrieben wird: zur Musik. Im aufgewühlten Paris der Februarrevolution beginnt die fulminante Suche nach dem tödlichen Geheimnis hinter der sagenhaften Seelenmacht von Frédéric Chopin.

      \medskip

      \begin{quote}
        \emph{%
‚Mein Lied willst du? Ich glaub es dir. Du wärst der erste Klavierspieler der Welt. Liszt, Chopin, Kalkbrenner würden verschwinden neben dir. Du würdest die Seelen bezwingen wie das Steppenfeuer dürres Gras. Auf den Knien würde die Welt vor dir liegen. Der Zar auf seinem Throne, Metternich, die Giftspinne der Heiligen Allianz […], sie beide würden zittern, würden sie dich dieses Lied spielen hören, gleichwie damals Ludwig XV. zitterte, als ich es ihm in Marly vorspielte.‘
        }
      \end{quote}
      \bigskip
    }}

    \bookcovercomponent{center rotate}{spine}{
      \color{white}\Large
      \bookauthor{} \hspace{8.4cm} \textbf{\booktitle{}}
    }
  \end{bookcover}
\end{document}
